% =====================================================
% 题型：平面向量平行、垂直与模长综合解答题 (q_vector_parallel_perpendicular_norm.tex)
% =====================================================
% =====================================================
% 难度：★★ (中档)
% =====================================================
\newcommand{\qVectorParallelPerpendicularNormStem}{已知平面向量 \(\vec{a} = (2, 1)\)，\(\vec{b} = (m-2, 4-2m)\)，\(\vec{c} = \vec{a} + \vec{b}\)。\\
（1）若 \(\vec{a} \parallel \vec{b}\)，求实数 \(m\) 的值；\\
（2）证明：对任意的 \(\lambda \ne 0\)，都有 \(\vec{a} \perp \lambda(\vec{a} - \vec{c})\)；\\
（3）若 \(\vec{a}\) 与 \(\vec{c}\) 的夹角为 \(\dfrac{\pi}{4}\)，求 \(|\vec{c}|\)。}

\newcommand{\qVectorParallelPerpendicularNormAnswer}{（1）\(m = 2\)；（2）见解析；（3）\(|\vec{c}| = \sqrt{10}\)}

\newcommand{\qVectorParallelPerpendicularNormSolution}{%
  【标准解答】
  
  （1）因为 \(\vec{a} = (2, 1)\)，\(\vec{b} = (m-2, 4-2m)\)，且 \(\vec{a} \parallel \vec{b}\)，
  根据向量平行的坐标表示，两向量的交叉乘积为 \(0\)：
  \[
    2(4-2m) - 1(m-2) = 0.
  \]
  展开并整理得：
  \[
    8 - 4m - m + 2 = 0 \implies 10 - 5m = 0 \implies m = 2.
  \]
  所以，当 \(\vec{a} \parallel \vec{b}\) 时，实数 \(m\) 的值为 \(2\)。
  
  \vspace{0.5em}
  
  （2）证明：因为 \(\vec{c} = \vec{a} + \vec{b}\)，所以 \(\vec{a} - \vec{c} = -\vec{b}\)。
  进而得：
  \[
    \lambda(\vec{a} - \vec{c}) = -\lambda \vec{b}.
  \]
  要证明 \(\vec{a} \perp \lambda(\vec{a} - \vec{c})\)，只需计算两者的数量积：
  \[
    \vec{a} \cdot \left[ \lambda(\vec{a} - \vec{c}) \right] = \vec{a} \cdot (-\lambda \vec{b}) = -\lambda (\vec{a} \cdot \vec{b}).
  \]
  计算 \(\vec{a} \cdot \vec{b}\) 的坐标点积：
  \[
    \vec{a} \cdot \vec{b} = 2(m-2) + 1(4-2m) = 2m - 4 + 4 - 2m = 0.
  \]
  由于 \(\vec{a} \cdot \vec{b} = 0\)，可知无论 \(m\) 为何值，都有 \(\vec{a} \perp \vec{b}\)。
  因此：
  \[
    \vec{a} \cdot \left[ \lambda(\vec{a} - \vec{c}) \right] = -\lambda \cdot 0 = 0.
  \]
  因为 \(\lambda \ne 0\)，数乘运算不改变垂直关系，所以对任意的 \(\lambda \ne 0\)，都有 \(\vec{a} \perp \lambda(\vec{a} - \vec{c})\)。
  
  \vspace{0.5em}
  
  （3）由（2）可知 \(\vec{a} \perp \vec{b}\)，即 \(\vec{a} \cdot \vec{b} = 0\)。
  因为 \(\vec{c} = \vec{a} + \vec{b}\)，我们计算 \(\vec{a} \cdot \vec{c}\)：
  \[
    \vec{a} \cdot \vec{c} = \vec{a} \cdot (\vec{a} + \vec{b}) = |\vec{a}|^2 + \vec{a} \cdot \vec{b} = |\vec{a}|^2.
  \]
  因为 \(\vec{a} = (2, 1)\)，所以：
  \[
    |\vec{a}| = \sqrt{2^2 + 1^2} = \sqrt{5}, \quad |\vec{a}|^2 = 5.
  \]
  因此，\(\vec{a} \cdot \vec{c} = 5\)。
  
  根据向量数量积的定义式，设 \(\vec{a}\) 与 \(\vec{c}\) 的夹角为 \(\theta = \dfrac{\pi}{4}\)，则：
  \[
    \vec{a} \cdot \vec{c} = |\vec{a}| |\vec{c}| \cos\theta = \sqrt{5} \cdot |\vec{c}| \cdot \cos\frac{\pi}{4} = \sqrt{5} \cdot |\vec{c}| \cdot \frac{\sqrt{2}}{2}.
  \]
  联立上述两式：
  \[
    5 = \sqrt{5} |\vec{c}| \cdot \frac{\sqrt{2}}{2} \implies |\vec{c}| = \frac{10}{\sqrt{10}} = \sqrt{10}.
  \]
  所以，\(|\vec{c}|\) 的值为 \(\sqrt{10}\)。
}

\newcommand{\qVectorParallelPerpendicularNormAltSolution}{%
  【几何多视角分析（直角三角形投影法）】
  
  由于对任意的实数 \(m\)，都有：
  \[
    \vec{a} \cdot \vec{b} = 2(m-2) + (4-2m) = 0,
  \]
  说明向量 \(\vec{a}\) 与 \(\vec{b}\) 始终保持垂直。
  
  在平面向量的加法几何意义中，若 \(\vec{a} \perp \vec{b}\)，且 \(\vec{c} = \vec{a} + \vec{b}\)，则以 \(\vec{a}, \vec{b}\) 为邻边可构造一个直角三角形，其中 \(\vec{c}\) 是斜边。
  
  根据直角三角形的性质，直角边 \(\vec{a}\) 在斜边 \(\vec{c}\) 上的投影长度为：
  \[
    |\vec{a}| = |\vec{c}| \cos\theta,
  \]
  其中 \(\theta\) 是 \(\vec{a}\) 与 \(\vec{c}\) 的夹角。
  
  当夹角 \(\theta = \dfrac{\pi}{4}\) 时，该直角三角形为等腰直角三角形。
  直接利用特殊角关系：
  \[
    |\vec{c}| = \frac{|\vec{a}|}{\cos\dfrac{\pi}{4}} = \sqrt{2} |\vec{a}|.
  \]
  因为 \(\vec{a} = (2,1) \implies |\vec{a}| = \sqrt{5}\)，
  所以：
  \[
    |\vec{c}| = \sqrt{2} \cdot \sqrt{5} = \sqrt{10}.
  \]
}

% =====================================================
% 别名兼容定义
% =====================================================
\newcommand{\qTJWVectorParallelPerpendicularNormStem}{\qVectorParallelPerpendicularNormStem}
\newcommand{\qTJWVectorParallelPerpendicularNormAnswer}{\qVectorParallelPerpendicularNormAnswer}
\newcommand{\qTJWVectorParallelPerpendicularNormSolution}{\qVectorParallelPerpendicularNormSolution}
\newcommand{\qTJWVectorParallelPerpendicularNormAltSolution}{\qVectorParallelPerpendicularNormAltSolution}
